\chapter{The reader's problem}
\label{ch:decision}

The worked note leaves an expert reader spending attention before the reader
knows which decision it supports. The funder's question is what a small,
credible intervention would have to do about that cost. Buying a writing bot,
certifying a document as ``human,'' and deploying an automated editor would be
different projects with different evidence.

\section{The problem worth paying to solve}

An expert reader's time is often the scarce input in a research proposal, a
model note, or an investment memorandum. The document may be mathematically
correct and professionally typeset while still asking that reader to infer why
the work matters, which evidence supports it, and what decision the author
wants. When the reader cannot do that, the review does not end. It returns to
the author as another round of diagnosis, revision, and rereading.

That failure occurred in a funding proposal in the project's internal record.
The source compiled. Automated checks found no fatal problem. Several rounds
of agent review made the prose smoother. A human reader still could not answer
three ordinary questions: What is this proposal trying to change? Why is the
proposed mechanism necessary? What action is being requested? The private file
names and phase labels that were meaningful to the authors occupied the space
where the argument should have been. The cost was not a bad sentence. It was a
transfer of reasoning from the author to a scarce decision maker.

This case is useful because it separates two things that are often confused. A
document can be locally clean and globally unsuccessful. A grammar
checker can be right about a comma while the reader is right that the document
has no visible point. A language model can produce a fluent paragraph while
changing the strength of a claim. Version control can show that a file changed
while saying nothing about whether the change helped.

\begin{takeawaybox}{The decision in one sentence}
Fund a bounded test of whether a revision workflow can reduce human-feedback
rounds without increasing reader time or changing meaning-bearing content.
\end{takeawaybox}

\section{What Humanvoice would do}

An author using Humanvoice would join three activities that are usually
scattered across tools and habits.

First, it locates problems that can be shown in the source. A finding points to
a passage, names the observable reason for attention, and asks an editorial
question when the judgment is substantive. The system can identify an internal
project identifier that has leaked into prose, repeated paragraph shapes,
defensive qualifications, a citation whose identity is inconsistent, or a
protected number that changed. It does not turn any of these observations into
a scalar quality score.

Second, it protects the objects that carry meaning while an author revises.
Equations, numbers, labels, citation keys, quotations, tables, code, and
declared claims receive source locations and a before-and-after comparison. An
author may change one of them, but the change cannot disappear inside a fluent
rewrite. The comparison makes revision inspectable rather than merely
trustworthy by reputation.

Third, it ends with a reader decision. A named reader sees the rendered
document without the repository history that made its private language seem
normal. The reader records whether the purpose, mechanism, evidence, and
requested action can be recovered. A model or metric can direct attention; a
human reader remains the authority who accepts a version or asks for a further
change.

The proposal is narrower than ``better writing.'' It asks whether connecting
located diagnosis, protected revision, and a reader outcome reduces wasted
review. Even a negative result is useful if it reveals which link creates
burden rather than saving it.

\section{The tools people already use}

The current alternative is a bundle of practices:

\begin{table}[H]
\centering
\small
\begin{tabularx}{\textwidth}{@{}p{0.23\textwidth}LL@{}}
\toprule
Practice & What it does well & What remains for the author and reader \\
\midrule
Prose or grammar linter & Finds local violations against explicit rules. & It
does not know which qualification, mechanism, or requested action matters. \\
General language model & Drafts, rewrites, and comments quickly. & It may
change a claim, number, or citation, and its preference is not reader
acceptance. \\
Version control and visual diff & Records that text changed and supports
rollback. & It does not identify which substantive change needs explanation. \\
Citation service & Resolves identifiers and metadata. & A valid DOI does not
show that the cited source supports the sentence. \\
Manual expert review & Can judge purpose, argument, and significance. & It is
scarce, repeated, and often receives a draft with no useful diagnostic trail. \\
\bottomrule
\end{tabularx}
\caption{The incumbent bundle and the gap Humanvoice would have to fill.}
\label{tab:incumbent-bundle}
\end{table}

Humanvoice should reuse these components where they work. Its reason for
existing is the transition between them: a finding should lead to an author
decision, a revision should lead to a protected comparison, and the comparison
should lead to a reader test. If the bundle can achieve the same result with
less cost, the product should not be built.

\section{The economic mechanism}

The relevant benefit is not a higher detector score or a prettier paragraph.
It is a reduction in the expert attention consumed before a document reaches a
sound decision. Let $N$ be the annual number of documents in scope, $r_0$ and
$r_1$ the feedback rounds under the incumbent and proposed workflows, $t_0$
and $t_1$ the expert-reader hours per round, $c_r$ the loaded cost of that
time, and $C$ the annualized product cost. The directly measurable component of
annual net value is
\begin{equation}
  V = N\bigl(r_0t_0-r_1t_1\bigr)c_r-C.
  \label{eq:value}
\end{equation}

The expression forces an important discipline. Saving an author ten minutes
does not create value if the reader spends another hour resolving false
alarms. A lower count of feedback rounds does not create value if the remaining
version changes a qualification or a number. The pilot must measure reader
time and meaning errors beside the round count.

The repository contains no defensible values for all five inputs, so the table
below is an illustration of the decision surface, not a forecast. It assumes
100 documents per year, a loaded reader cost of 150 currency units per hour,
and an annual product cost of 60,000 units. The rows vary the hours saved per
document; they show why baseline measurement matters more than a confident
point estimate.

\begin{table}[H]
\centering
\small
\begin{tabular}{@{}lrrr@{}}
\toprule
Reader hours saved per document & Annual reader hours saved & Gross value & Net value \\
\midrule
1 & 100 & 15,000 & $-45,000$ \\
4 & 400 & 60,000 & 0 \\
8 & 800 & 120,000 & 60,000 \\
\bottomrule
\end{tabular}
\caption{Illustrative sensitivity, not a Humanvoice estimate. The feasibility
run supplies the inputs that determine which row is relevant.}
\label{tab:value-sensitivity}
\end{table}

The same calculation can be expressed in terms of feedback rounds when reader
time per round is stable. That is useful for the pilot, but it is not a license
to treat rounds as the whole outcome. A review round is a process measure; the
decision maker cares whether the final document is understood and whether its
meaning survived.

\section{What the evidence permits us to say now}

The evidence supports a bounded investment, not a deployment claim. Controlled
experiments show that generative assistance can make some writing and
knowledge-work tasks faster or better scored, while field results vary with
task and worker experience \citep{noy2023experimental,brynjolfsson2025work,dillon2025shifting}.
Work outside a model's capability frontier can reduce correctness even when
the output appears useful \citep{dellacqua2026jagged}. Studies of writing
feedback find that models can identify local defects while missing the central
problem in a passage \citep{nair2024closing,rashkin2025story}. Recent reviews
find more reliable gains in grammar, conventions, and process than in
argumentation or discourse organization \citep{meng2026systematic,abudalfa2026awe}.

Those findings make the proposed division of labor sensible: deterministic
tools locate what they can observe, an author decides what a repair means, and
a reader supplies the acceptance outcome. They do not show that the combined
workflow reduces revision cost for expert economics writing. That is the
question the feasibility build must prepare us to test.

\section{What the first build will and will not do}

The request is deliberately narrow: implement the local-first
document core and high-value diagnostics; assemble a small corpus with source
snapshots, protected objects, and human labels; run one live document through
the complete path; and return with a priced proceed, revise, or stop
recommendation. It is not permission to deploy, classify authorship,
optimize against an AI detector, or claim improved expert writing. The product
must first show that its observed workflow earns its cost, protects meaning,
and leaves a named reader able to make the intended decision.
